\section*{Bab \ref{ch:b1:curves} --- Kurva Bidang}

\begin{solution}{exo:b1:curves:1}
$M(-t) = (t^2, -t^3)$: yaitu pencerminan pada sumbu $x$; jadi pelajarilah $t \geq 0$.
Baik $x' = 2t$ maupun $y' = 3t^2$ bernilai $\geq 0$: sehingga cabangnya bergerak
ke kanan dan ke atas, dari $(0,0)$ ke takhingga. Di $t = 0$ kecepatannya
lenyap; dan penjabaran $x = t^2$, $y = t^3$ menunjukkan bahwa garis singgungnya
sumbu $x$ ($y/x = t \to 0$) dengan $y$ berganti tanda sedangkan $x \geq 0$:
jadi sebuah \emph{\omterm{ex:b1:curves:astroid}{titik runcing}} yang menunjuk ke kiri. Kurvanya adalah parabola semikubik
$y^2 = x^3$.
\end{solution}

\begin{solution}{exo:b1:curves:2}
$M(t + 2\pi) = M(t) + (2\pi, 0)$: jadi kurvanya berulang, tergeser sejauh
satu keliling roda; pelajarilah satu periode. $x'(t) = 1 - \cos t \geq
0$, $y'(t) = \sin t$: sehingga lengkungannya naik pada $\intoo{0}{\pi}$, turun pada
$\intoo{\pi}{2\pi}$, dan memuncak di $(\pi, 2)$. Adapun titik singularnya
tempat $x' = y' = 0$: yaitu $t \in 2\pi\Z$, pada tanahnya. Dekat $t = 0$:
\[
x(t) = \frac{t^3}{6} + o(t^3),
\qquad
y(t) = \frac{t^2}{2} + o(t^2):
\]
$x$ berganti tanda, $y \geq 0$, dan $\frac{x}{y^{3/2}}$ terbatas: sehingga
garis singgungnya \emph{tegak} (berarah $(0,1)$), yaitu \omterm{ex:b1:curves:astroid}{titik runcing} tempat
titik yang dijejaki sesaat berlaju nol --- yakni tanda fisik
menggelinding tanpa tergelincir.
\end{solution}

\begin{solution}{exo:b1:curves:3}
Di $t = \frac\pi4$: $M = \bigl(\frac{\sqrt2}{4},
\frac{\sqrt2}{4}\bigr)$, dengan arah singgung $(-\cos t, \sin t) =
\frac{1}{\sqrt2}(-1, 1)$ (\cref{ex:b1:curves:astroid}). Garis
singgungnya: $y - \frac{\sqrt2}{4} = -(x - \frac{\sqrt2}{4})$, yaitu $x +
y = \frac{\sqrt2}{2}$. Perpotongannya: $\bigl(\frac{\sqrt2}{2},
0\bigr)$ dan $\bigl(0, \frac{\sqrt2}{2}\bigr)$; sehingga ruas di antara
keduanya berpanjang $\sqrt{\frac12 + \frac12} = 1$.

(Titik umum $t$: garis singgung di $(\cos^3 t, \sin^3 t)$ memotong
sumbunya di $(\cos t, 0)$ dan $(0, \sin t)$ --- periksalah bahwa garis
lewat titik ini berarah $(-\cos t, \sin t)$ dan lewat
$M(t)$ --- dan ruas yang terpotong berpanjang $\sqrt{\cos^2 t +
\sin^2 t} = 1$.)
\end{solution}

\begin{solution}{exo:b1:curves:4}
$r = \cos 2\theta$: berperiode $\pi$ pada $\theta$, setangkup pada kedua
sumbunya; $r$ lenyap di $\theta = \pm\frac\pi4$ (dengan garis singgung di
titik asalnya sepanjang diagonalnya) dan negatif untuk $\theta \in
\intoo{\frac\pi4}{\frac{3\pi}{4}}$, tempat titiknya terplot pada
sinar garis yang berlawanan --- sehingga menghasilkan empat kelopak sepanjang arah sumbunya
$\theta = 0, \frac\pi2, \pi, \frac{3\pi}{2}$, yang masing-masingnya berjari-jari
maksimal $1$.

$r\cos\theta = 1$ adalah persamaan $x = 1$: sehingga \omterm{def:b1:curves:polar}{kurva kutub} $r =
\frac{1}{\cos\theta}$ merupakan garis tegak $x = 1$ (yang tersapu sekali untuk
$\theta \in \intoo{-\frac\pi2}{\frac\pi2}$).
\end{solution}

\begin{solution}{exo:b1:curves:5}
$r(\theta) = 1 + 2\cos\theta$ lenyap untuk $\cos\theta = -\frac12$:
yaitu $\theta = \pm\frac{2\pi}{3}$. Kesetangkupannya pada sumbu $x$; jadi pelajarilah
$\theta \in \intcc{0}{\pi}$. Untuk $\theta \in
\intco{0}{\frac{2\pi}{3}}$: $r > 0$, menurun dari $3$ ke $0$: yaitu
busur luarnya, yang masuk ke titik asalnya menyinggung sinar garis $\theta =
\frac{2\pi}{3}$. Untuk $\theta \in \intoc{\frac{2\pi}{3}}{\pi}$: $r <
0$: sehingga titik $r\vec u(\theta)$ duduk pada sinar garis yang berlawanan (bersudut
$\theta - \pi \in \intoc{-\frac\pi3}{0}$), dengan jarak $\abs r$
tumbuh dari $0$ ke $1$: dan ini menggambar sebuah \emph{simpul dalam} yang kecil
lewat titik asalnya. Di $\theta = \pi$, $r = -1$ dan titiknya adalah
$-\vec u(\pi) = (1, 0)$: jadi simpulnya menutup pada sumbu $x$ yang positif.
Sketsanya: sebuah kurva luar yang besar bagai hati
dengan jangkauan maksimal $3$ di $\theta = 0$, ditambah sebuah simpul di dalamnya
lewat $O$ dan $(1,0)$.
\end{solution}

\begin{solution}{exo:b1:curves:6}
Lengkapkanlah kuadratnya:
\[
9(x^2 - 4x) + 25(y^2 - 2y) = 164
\iff 9(x-2)^2 + 25(y-1)^2 = 164 + 36 + 25 = 225 .
\]
\omterm{def:b1:arith:divides}{Membaginya} dengan $225$: $\frac{(x-2)^2}{25} + \frac{(y-1)^2}{9} = 1$: yaitu
elips berpusat $(2, 1)$, dengan setengah sumbu $a = 5$ (mendatar), $b =
3$; $c = \sqrt{25 - 9} = 4$: fokusnya $(2 \pm 4,\, 1) = (-2, 1)$ dan
$(6, 1)$; dan eksentrisitasnya $e = \frac45$.
\end{solution}

\begin{solution}{exo:b1:curves:7}
Fokusnya di titik asalnya, direktriksnya $D: x = d$ dengan $d = \frac pe > 0$.
Bagi sebuah titik $M = (r\cos\theta, r\sin\theta)$ dengan $r > 0$:
\[
MF = r,
\qquad
d(M, D) = \abs{d - r\cos\theta},
\]
dan syarat \omterm{def:b1:curves:conics}{irisan kerucutnya} $MF = e\,d(M, D)$, pada rezim tempat $M$
ada di sisi fokus direktriksnya ($r\cos\theta < d$), terbaca $r
= e(d - r\cos\theta)$, yakni
\[
r(1 + e\cos\theta) = ed = p,
\qquad
r = \frac{p}{1 + e\cos\theta} .
\]
Untuk $e < 1$ penyebutnya tak pernah lenyap: sehingga semua $\theta$ diperkenankan
(yaitu elips). Untuk $e = 1$: $\theta \neq \pi$ (yaitu parabola, yang \omterm{def:b1:topology:open}{terbuka} ke arah
sisi jauh direktriksnya). Untuk $e > 1$: dituntut $1 + e\cos\theta > 0$,
yaitu $\theta \in \intoo{-\theta_0}{\theta_0}$ dengan $\theta_0 =
\arccos\bigl(-\frac1e\bigr)$: jadi satu \emph{cabang} hiperbolanya
(sedangkan cabang yang lain bersesuaian dengan pilihan tanda $r < 0$, atau dengan
fokus keduanya).
\end{solution}

\begin{solution}{exo:b1:curves:8}
\emph{Tukang kebun:} ikatkanlah seutas tali berpanjang $2a$ pada dua pasak $F,
F'$ lalu tegangkanlah ia dengan titik penjejaknya $M$: maka kendalanya tepat
$MF + MF' = 2a$, sehingga kurva yang terjejaki adalah elipsnya.

\emph{Sifat pemantulannya:} misalkan $t \mapsto M(t)$ sebuah pemarameteran
yang reguler dan $u(t) = \frac{M(t) - F}{\norm{M(t) - F}}$,
dengan $u'(t)$ vektor satuan yang serupa menuju $F'$. Menurunkan
$\norm{M - F} + \norm{M - F'} = 2a$, dengan memakai $\frac{\dd}{\dd
t}\norm{M - F} = \bigl\langle \frac{M - F}{\norm{M-F}},\,
M'\bigr\rangle$ (yaitu kaidah rantai pada $\sqrt{\langle\cdot,\cdot\rangle}$):
\[
\bigl\langle u + u',\, M'\bigr\rangle = 0 .
\]
Jadi arah singgungnya $M'$ \omterm{def:b1:euclid:orthogonal}{ortogonal} terhadap arah garis bagi
$u + u'$ atas kedua sinar garis fokusnya: sehingga garis singgungnya bersudut sama
dengan $MF$ dan $MF'$. (Karena seberkas cahaya dari satu fokusnya memantul
pada elipsnya menuju fokus yang lain.)
\end{solution}

\begin{solution}{exo:b1:curves:9}
\begin{enumerate}
  \item Daerah asalnya: $\cos 2\theta \geq 0$: yaitu $\theta \in
        \intcc{-\frac\pi4}{\frac\pi4} \cup
        \intcc{\frac{3\pi}{4}}{\frac{5\pi}{4}}$. Kesetangkupannya: $\theta
        \mapsto -\theta$ (sumbu $x$) dan $\theta \mapsto \pi -
        \theta$ (sumbu $y$): jadi pelajarilah $\intcc{0}{\frac\pi4}$; dan $r$
        menurun dari $1$ ke $0$. Di $\theta = \pm\frac\pi4$: $r =
        0$, dengan garis singgung di titik asalnya sepanjang diagonalnya. Kurvanya adalah
        lambang $\infty$: yaitu dua simpul setangkup yang bertemu di $O$,
        dan mencapai $(\pm 1, 0)$.
  \item Dengan $F, F' = (\pm c, 0)$, $c = \frac{1}{\sqrt 2}$, dan $M
        = (r\cos\theta, r\sin\theta)$:
        \[
        MF^2\, MF'^2
        = \bigl((r\cos\theta - c)^2 + r^2\sin^2\theta\bigr)
        \bigl((r\cos\theta + c)^2 + r^2\sin^2\theta\bigr)
        = (r^2 + c^2)^2 - (2rc\cos\theta)^2 ,
        \]
        lewat kesamaan $(A - B)(A + B) = A^2 - B^2$ dengan $A = r^2
        + c^2$, $B = 2rc\cos\theta$. Dengan $c^2 = \frac12$:
        \[
        MF^2 MF'^2 = \Bigl(r^2 + \frac12\Bigr)^2 - 2r^2\cos^2\theta
        = r^4 + r^2\bigl(1 - 2\cos^2\theta\bigr) + \frac14
        = r^4 - r^2\cos 2\theta + \frac14 .
        \]
        Pada lemniskatnya, $r^2 = \cos 2\theta$: sehingga kedua suku pertamanya
        saling menghapus, dan menyisakan $MF^2MF'^2 = \frac14$, yaitu $MF \cdot MF' =
        \frac12$. Sebaliknya, perhitungan yang dibaca terbalik menunjukkan
        bahwa persamaan tempat kedudukannya $MF\,MF' = \frac12$ adalah $r^2 = \cos
        2\theta$ (untuk $r \neq 0$; dan $O$ memenuhi keduanya).
\end{enumerate}
\end{solution}

\begin{solution}{exo:b1:curves:10}
$r = 1 + \cos\theta$, $r' = -\sin\theta$:
\[
r^2 + r'^2 = 1 + 2\cos\theta + \cos^2\theta + \sin^2\theta
= 2 + 2\cos\theta = 4\cos^2\frac\theta2 ,
\]
sehingga $\sqrt{r^2 + r'^2} = 2\abs{\cos\frac\theta2}$ dan, lewat
kesetangkupannya pada sumbu $x$,
\[
L = \int_0^{2\pi} 2\Bigl|\cos\frac\theta2\Bigr|\,\dd\theta
= 2\int_0^{\pi} 2\cos\frac\theta2\,\dd\theta
= 2\Bigl[4\sin\frac\theta2\Bigr]_0^{\pi} = 8 .
\]
Sekali lagi sebuah keliling yang aljabar dan \omterm{def:b1:vspaces:free}{bebas} $\pi$.
\end{solution}

\begin{solution}{exo:b1:curves:11}
Titik $(a\cos t, b\sin t)$ memenuhi persamaannya:
$\cos^2 t + \sin^2 t = 1$. Adapun vektor normal garisnya adalah
$\bigl(\frac{\cos t}a, \frac{\sin t}b\bigr)$, dan hasil kalinya
dengan kecepatannya $(-a\sin t, b\cos t)$ adalah $-\sin t\cos t +
\sin t\cos t = 0$: sehingga garisnya lewat titiknya dengan
arah singgungnya --- jadi ialah garis singgungnya. Bagi $(x_0, y_0)$ pada
elipsnya, tulislah $\cos t = \frac{x_0}a$, $\sin t = \frac{y_0}b$ lalu
sulihkanlah:
\[
\frac{x\,x_0}{a^2} + \frac{y\,y_0}{b^2} = 1 ,
\]
yang diperoleh dari persamaan elipsnya dengan ``membelah'' $x^2 \mapsto
x\,x_0$ dan $y^2 \mapsto y\,y_0$.
\end{solution}

\begin{solution}{exo:b1:curves:12}
\begin{enumerate}
  \item $r' = k\eu^{k\theta}$, sehingga $\tan V = \frac{r}{r'} =
        \frac1k$: jadi tetap. Sehingga spiralnya memotong setiap sinar garis dari
        titik asalnya pada sudut yang sama $V = \arctan\frac1k$.
  \item Dalam notasi kompleks spiralnya adalah $\{\eu^{k\theta}
        \eu^{\iu\theta} : \theta \in \R\}$. Memutarnya sebesar $c$
        mengalikannya dengan $\eu^{\iu c}$:
        \[
        \eu^{k\theta}\eu^{\iu(\theta + c)}
        = \eu^{-kc}\,\eu^{k(\theta + c)}\eu^{\iu(\theta+c)} ,
        \]
        dan ketika $\theta + c$ menjelajahi $\R$ ini memerikan
        $\eu^{-kc}$ kali spiralnya: jadi rotasi $=$ penskalaan. Tak ada
        kurva mulus lain selain garis dan lingkaran yang mempunyai
        sifat ini.
  \item $\sqrt{r^2 + r'^2} = \sqrt{1 + k^2}\,\eu^{k\theta}$, sehingga
        pada $\intcc{A}{\theta_0}$ \omterm{def:b1:curves:arclength}{panjangnya} adalah
        $\frac{\sqrt{1+k^2}}{k}\bigl(\eu^{k\theta_0} -
        \eu^{kA}\bigr)$, dan ketika $A \to -\infty$:
        \[
        L = \frac{\sqrt{1 + k^2}}{k}\,\eu^{k\theta_0} < \infty :
        \]
        yaitu takhingga banyaknya putaran mengelilingi titik asalnya, tetapi berpanjang
        total hingga (karena putarannya menyusut secara geometri).
\end{enumerate}
\end{solution}

\begin{solution}{pb:b1:curves:1}
\textbf{1.} Menggelinding tanpa tergelincir berarti titik sentuhnya telah
menempuh jarak yang sama dengan busur roda yang tergulung: sehingga setelah
berputar sebesar $t$, pusatnya ada di $(Rt, R)$. Adapun titik yang ditandai duduk
pada peleknya bersudut $t$ \emph{di belakang} garis tegak ke bawahnya
(karena rodanya berputar searah jarum jam sembari maju):
\[
M(t) = \Omega(t) + R(-\sin t, -\cos t)
= \bigl(R(t - \sin t),\ R(1 - \cos t)\bigr),
\]
yang benar di $t = 0$ ($M = (0,0)$) dan di $t = \pi$ ($M =
(\pi R, 2R)$, yaitu puncaknya).

\textbf{2.} $f'(t) = R(1 - \cos t,\ \sin t)$ dan
\[
\norm{f'}^2 = R^2\bigl((1-\cos t)^2 + \sin^2 t\bigr)
= 2R^2(1 - \cos t) = 4R^2\sin^2\frac t2 ,
\]
sehingga $\norm{f'} = 2R\abs{\sin\frac t2}$. Titik singularnya di $t \in
2\pi\Z$: yaitu \omterm{ex:b1:curves:astroid}{titik runcing} pada tanahnya (\cref{exo:b1:curves:2}); sedangkan
puncak lengkungannya di $t = \pi$, tempat lajunya $2R$ maksimal.

\textbf{3.} Pemfaktoran setengah sudutnya:
\[
f'(t) = 2R\sin\frac t2\,\Bigl(\sin\frac t2,\ \cos\frac t2\Bigr),
\quad
M - C = 2R\sin\frac t2\,\Bigl(-\cos\frac t2,\ \sin\frac t2\Bigr),
\]
\[
T - M = \bigl(R\sin t,\ R(1 + \cos t)\bigr)
= 2R\cos\frac t2\,\Bigl(\sin\frac t2,\ \cos\frac t2\Bigr).
\]
Untuk $0 < t < 2\pi$, $\sin\frac t2 \neq 0$: sehingga $M - C$ \omterm{def:b1:euclid:orthogonal}{ortogonal}
terhadap arah singgungnya $\bigl(\sin\frac t2, \cos\frac
t2\bigr)$ (karena hasil kalinya adalah $-\sin\frac t2\cos\frac t2 +
\sin\frac t2\cos\frac t2 = 0$), sedangkan $T - M$ sejajar dengannya.
Jadi tali busur ke puncak rodanya \emph{adalah} garis singgungnya; sedangkan
tali busur ke titik sentuhnya adalah normalnya.

\textbf{4.} Pada setiap saat rodanya berporos pada titik
sentuhnya (karena titik itu berkecepatan nol: yakni menggelinding tanpa tergelincir),
sehingga setiap titik tegar rodanya bergerak \omterm{def:b1:euclid:orthogonal}{ortogonal} terhadap garis
yang menghubungkannya ke $C$, dengan laju (pada kecepatan sudut $1$) yang sama dengan
jarak itu. Periksa: $\norm{M - C} = 2R\abs{\sin\frac t2} =
\norm{f'(t)}$.

\textbf{5.} $2R\,y(t) = 2R\cdot 2R\sin^2\frac t2 =
4R^2\sin^2\frac t2 = \norm{f'(t)}^2$. Jadi lajunya pada ketinggian $y$ adalah
$\sqrt{2R\,y}$ --- yang secara bentuk adalah hukum $v = \sqrt{2gh}$ pada jatuh
\omterm{def:b1:vspaces:free}{bebas}, dengan tanahnya berperan sebagai langit-langitnya; dan Bagian IV mengubah
pengamatan ini menjadi kerja jam.

\textbf{6.} Menurut \cref{def:b1:curves:arclength} dan pertanyaan 2
($\sin\frac t2 \geq 0$ pada $\intcc{0}{2\pi}$):
\[
L = \int_0^{2\pi} 2R\sin\frac t2\,\dd t
= 2R\Bigl[-2\cos\frac t2\Bigr]_0^{2\pi} = 2R(2 + 2) = 8R .
\]
Teorema Wren: tepat empat garis tengah.

\textbf{7.} $s(t) = \int_0^t 2R\sin\frac u2\,\dd u =
4R\bigl(1 - \cos\frac t2\bigr)$; dan $s(2\pi) = 4R(1+1) = 8R$,
yang selaras.

\textbf{8.} $\sigma(t) = \abs{s(t) - s(\pi)} = \abs{4R(1 -
\cos\frac t2) - 4R} = 4R\abs{\cos\frac t2}$, sehingga $\sigma^2 =
16R^2\cos^2\frac t2$; dan $2R - y = R(1 + \cos t) =
2R\cos^2\frac t2$, dan dari situlah $8R(2R - y) = 16R^2\cos^2\frac t2 =
\sigma^2$.

\textbf{9.} Pada \omterm{ex:b1:curves:astroid}{titik runcing} $t = 0$ (atau $2\pi$): $\sigma = 4R$, yaitu separuh
$8R$: jadi puncaknya membelah dua lengkungannya. Pada kedua perhitungannya
lajunya berupa $\abs{\text{polinomial trigonometri pada } t/2}$, yang
antiturunannya trigonometri lagi: sehingga \omterm{def:b1:curves:arclength}{panjangnya} adalah
selisih \emph{nilai} kosinus --- yaitu bilangan rasional
kali $R$ --- tanpa busur lingkaran yang perlu diukur, jadi tanpa $\pi$.
Sedangkan lingkarannya sendiri berlaju tetap, sehingga \omterm{thm:b1:integration:def}{integral} \omterm{def:b1:curves:arclength}{panjangnya}
menghasilkan \omterm{def:b1:curves:arclength}{panjang} \omterm{prop:b1:reals:intervals}{selang} penuhnya $2\pi$; sedangkan laju sikloidnya
lenyap pada ujungnya dan terintegralkan secara aljabar.

\textbf{10.} Lengkungannya tersapu dengan $x$ naik dari $0$ ke
$2\pi R$ ($x' = R(1 - \cos t) \geq 0$, yang lenyap hanya pada
titik yang terpencil). Jadi menyulihkan $x = x(t)$ ke dalam \omterm{thm:b1:integration:def}{integral} luasnya
$\int_0^{2\pi R} y\,\dd x$ (\cref{thm:b1:integration:parts})
memberikan $A = \int_0^{2\pi} y(t)\,x'(t)\,\dd t$.

\textbf{11.}
\[
A = \int_0^{2\pi} R(1 - \cos t)\cdot R(1 - \cos t)\,\dd t
= R^2\int_0^{2\pi}\bigl(1 - 2\cos t + \cos^2 t\bigr)\dd t
= R^2\Bigl(2\pi - 0 + \pi\Bigr) = 3\pi R^2 ,
\]
dengan memakai $\int_0^{2\pi}\cos^2 = \pi$. Jadi tepat tiga luas roda:
karena timbangan Galileo berkata ``sekitar $3$''; sedangkan perhitungan Roberval
berkata ``tepat''.

\textbf{12.} Dengan $x = \cos^3 t$, $y = \sin^3 t$: $y\,x' =
-3\sin^4 t\cos^2 t$. Linearkanlah:
\[
\sin^4 t\cos^2 t = \frac{1 - \cos 2t}{2}\cdot\frac{\sin^2
2t}{4} = \frac{\sin^2 2t}{8} - \frac{\sin^2 2t\cos 2t}{8} ,
\]
dan atas $\intcc{0}{2\pi}$: $\int\sin^2 2t = \pi$, $\int \sin^2
2t\cos 2t = \bigl[\frac{\sin^3 2t}{6}\bigr] = 0$. Jadi $\oint
y\,\dd x = -3\cdot\frac\pi8$, dan luas yang dilingkupinya adalah
$\frac{3\pi}8$ (dengan tandanya mencatat arah berlawanan
jarum jam).

\textbf{13.} Untuk $(\cos t, \sin t)$ dengan $t$ dari $\pi$ ke $0$,
$x$ naik dari $-1$ ke $1$ dan
\[
\int_\pi^0 \sin t\cdot(-\sin t)\,\dd t
= \int_0^\pi \sin^2 t\,\dd t = \frac\pi2 :
\]
jadi rumusnya mengembalikan luas setengah cakram atasnya, sebagaimana mestinya.

\textbf{14.} $x' = R(1 + \cos t) = 2R\cos^2\frac t2$, $y' =
R\sin t = 2R\sin\frac t2\cos\frac t2$, sehingga $\norm{f'} =
2R\cos\frac t2$ (yang taknegatif pada $\intcc{-\pi}{\pi}$). Adapun busur dari
dasarnya: $s(t) = \int_0^t 2R\cos\frac u2\,\dd u =
4R\sin\frac t2$. Lalu
\[
y = R(1 - \cos t) = 2R\sin^2\frac t2
= 2R\Bigl(\frac{s}{4R}\Bigr)^{\!2} = \frac{s^2}{8R} .
\]

\textbf{15.} Kekekalan energinya dengan $v = \frac{\dd s}{\dd\tau}$
dan $y = \frac{s^2}{8R}$, $y_0 = \frac{s_0^2}{8R}$:
\[
\Bigl(\frac{\dd s}{\dd\tau}\Bigr)^{\!2} = 2g(y_0 - y)
= \frac{2g}{8R}\bigl(s_0^2 - s^2\bigr)
= \frac{g}{4R}\bigl(s_0^2 - s^2\bigr).
\]

\textbf{16.} Menurunkannya pada $\tau$: $2s's'' =
-\frac{g}{4R}\,2ss'$, sehingga di mana pun $s' \neq 0$ (jadi di mana-mana
lewat \omterm{def:b1:continuity:continuous}{kekontinuannya}), $s'' = -\frac{g}{4R}s$: yaitu pengayun harmonisnya.
Menurut \cref{thm:b1:diffeq:homogeneous2}, $s(\tau) =
A\cos\omega\tau + B\sin\omega\tau$ dengan $\omega =
\sqrt{g/(4R)}$; dan syarat awalnya $s(0) = s_0$, $s'(0) = 0$
memberikan $s(\tau) = s_0\cos\omega\tau$.

\textbf{17.} Maniknya mencapai dasarnya ketika $s = 0$, yaitu di
$\omega\tau = \frac\pi2$:
\[
T_{\downarrow} = \frac{\pi}{2\omega}
= \frac\pi2\sqrt{\frac{4R}{g}} = \pi\sqrt{\frac Rg}\,,
\]
yang tak bergantung pada $s_0$: jadi dilepas dari mana pun di mangkuknya, semua
maniknya tiba pada saat yang sama --- itulah tautokronnya.

\textbf{18.} Menyulihkan $s = s_0\cos\omega\tau$: maka ruas kirinya
adalah $s_0^2\omega^2\sin^2\omega\tau$ sedangkan ruas kanannya adalah
$\frac{g}{4R}s_0^2(1 - \cos^2\omega\tau) =
s_0^2\omega^2\sin^2\omega\tau$: jadi kesamaan yang persis, sehingga tak ada
penyelesaian palsu yang diperkenalkan. Adapun bagi busur lingkaran, $y$ tak
sebanding dengan $s^2$ ($y = \ell(1 - \cos\frac s\ell) =
\frac{s^2}{2\ell} - \frac{s^4}{24\ell^3} + \dots$): sehingga
suku pemulihnya hanya \emph{secara hampiran} \omterm{def:b1:linmaps:def}{linear}, jadi
periode bandul lingkarannya melayang bersama amplitudonya, sedangkan
periode sikloidnya tetap secara ketat.

\textbf{19.} $\pi\sqrt{R/g} = 1$ memberikan $R = \frac{g}{\pi^2}
\approx \frac{9.81}{9.87} \approx 0.994$ m --- yaitu nyaris tepat
satu meter. Adapun ayunan penuhnya memakan $\frac{2\pi}\omega =
2\pi\sqrt{4R/g}$, yang persis rumus bersudut kecil
$2\pi\sqrt{\ell/g}$ bagi bandul berpanjang $\ell = 4R$:
dan Huygens menggantungkan bandulnya pada pipi sikloidal dengan
perbandingan tepat itu.

\textbf{20.} Di $t = \frac\pi2$, $R = 1$: $M = \bigl(\frac\pi2 -
1,\ 1\bigr)$, $T = \bigl(\frac\pi2,\ 2\bigr)$, sehingga $T - M = (1,
1)$. Dan $f'(\frac\pi2) = (1 - 0,\ 1) = (1, 1)$: jadi tali busur ke
puncaknya adalah kecepatannya, sebagaimana dijanjikan.

\textbf{21.} Bagi trokoidnya, $x'(t) = R - d\cos t$ dan $y'(t)
= d\sin t$. Jika $d < R$: maka $x' \geq R - d > 0$, sehingga titiknya selalu
maju; dan kecepatannya tak pernah lenyap (karena komponen pertamanya
positif): jadi tak ada titik singular, melainkan gelombang yang reguler. Jika $d > R$: $x'(0)
= R - d < 0 < R + d = x'(\pi)$, sehingga titiknya bergerak mundur
dekat saat sentuhnya dan maju di tempat lain: jadi kurvanya
menyilang dirinya dalam simpul. Adapun sebuah titik pada tepi roda
kereta api (di bawah kepala relnya, $d > R$) bepergian mundur pada setiap
putaran.

\textbf{22.} Tali busur dari \omterm{ex:b1:curves:astroid}{titik runcing} $(\pi R, 2R)$ ke titik asalnya
berpanjang $L = R\sqrt{\pi^2 + 4}$ dan bersudut kemiringan $\alpha$ dengan
$\sin\alpha = \frac{2R}{L}$. Meluncur dari diam dengan percepatan
tetap $g\sin\alpha$: $L = \frac12 g\sin\alpha\,T^2$, sehingga
\[
T = \sqrt{\frac{2L}{g\sin\alpha}} = \sqrt{\frac{2L^2}{2Rg}}
= \frac{L}{\sqrt{Rg}} = \sqrt{\pi^2 + 4}\,\sqrt{\frac Rg}
\approx 3.72\sqrt{\frac Rg}\,,
\]
lawan $\pi\sqrt{R/g} \approx 3.14\sqrt{R/g}$ bagi sikloidnya:
jadi lintasan yang melengkung lebih cepat. Dan ia sesungguhnya yang paling cepat
--- yakni brakistokronnya --- yaitu teorema kalkulus
variasi, yang di luar jilid ini.

\textbf{23.} $\dfrac{\dd y}{\dd x} = \dfrac{y'}{x'} =
\dfrac{\sin t}{1 - \cos t} = \cot\dfrac t2$ (lewat rumus
setengah sudutnya). Lalu
\[
\frac{\dd^2y}{\dd x^2}
= \frac{\dd}{\dd t}\Bigl(\cot\frac t2\Bigr)\cdot\frac{1}{x'(t)}
= -\frac{1}{2\sin^2\frac t2}\cdot\frac{1}{2R\sin^2\frac t2}
= -\frac{1}{4R\sin^4\frac t2} < 0 :
\]
jadi cekung sepanjang lengkungannya; dan ketika $t \to 0^+$ atau $2\pi^-$
kemiringannya $\cot\frac t2 \to \pm\infty$: jadi garis singgung tegak pada
\omterm{ex:b1:curves:astroid}{titik runcingnya}.

\textbf{24.} Arah tali busurnya adalah $T - M \parallel
\bigl(\sin\frac t2, \cos\frac t2\bigr)$, yang menuju $(0, 1)$
ketika $t \to 0^{+}$: sehingga garis singgung pada \omterm{ex:b1:curves:astroid}{titik runcingnya} tegak ---
yang dipulihkan dari geometri murni, tanpa penjabaran Taylor.

\textbf{25.} (i) Teorema Wren, $L = 8R$; luas Roberval, $A =
3\pi R^2$ (dengan nisbah $3$ yang diterka Galileo); dan tautokron
Huygens, $T_\downarrow = \pi\sqrt{R/g}$. (ii) Setiap
perhitungannya berjalan di atas pemfaktoran setengah sudutnya $f' =
2R\sin\frac t2(\sin\frac t2, \cos\frac t2)$ dan padanan mangkuknya
--- yaitu satu kesamaan yang menggerakkan garis singgung, \omterm{def:b1:curves:arclength}{panjang} dan jam
sekaligus. (iii) Adapun hubungan intrinsiknya $y = s^2/8R$ mengubah
persamaan energinya menjadi $s'' = -\frac{g}{4R}s$, yaitu persamaan \omterm{def:b1:linmaps:def}{linear}
berkoefisien tetap yang penyelesaiannya persis
isokron. (iv) Bagian I memakai kalkulus diferensial
\cref{ch:b1:derivative}, Bagian II--III \omterm{thm:b1:integration:def}{integral}
\cref{ch:b1:integration}, dan Bagian IV persamaan diferensial
\cref{ch:b1:diffeq} --- jadi sikloidnya adalah seluruh bahan ajar buku ini
yang tergulung menjadi satu kurva.

\end{solution}
